\noindent\textbf{Reviewer comment:}
\emph{Besides, the paper must extend simulation verification by incorporating
communication packet loss, sensor noise, and partial interceptor failures to
validate robustness in the proposed cooperative swarm interception framework.}

\medskip
\noindent\textbf{Response:}
Thank you for this important suggestion. We added a
partial-interceptor-failure study using the frozen policies and the same
three-dimensional evaluation environment. In each of 100 independent trials
per case, two of the 20 interceptors were sampled uniformly without
replacement and made unavailable after target assignment but before the first
guidance command. No retraining, online reassignment, or post-selection of
trials was performed. Mission success required all eight attackers to be
intercepted by the remaining operational vehicles. The resulting interception
success rates were 92.0\% in Case~1 and 44.0\% in Case~2. For successful
trials, the mean
$\{E_{co\text{-}time},E_n,E_{miss},E_t\}$ values were
$\{0.0113~\mathrm{s},0.2616~g,1.6582~\mathrm{m},33.7859~\mathrm{s}\}$
and
$\{0.8520~\mathrm{s},0.2628~g,2.0530~\mathrm{m},61.6489~\mathrm{s}\}$,
respectively.

Failure decomposition further clarifies the limitation. In each case, only
five failed trials were caused by both unavailable interceptors removing all
vehicles assigned to one target. In Case~2, 51 of the 56 failed trials
contained an uncovered target for which at least one operational interceptor
remained. Thus, the fixed assignment provides useful failure redundancy in
Case~1, whereas the more demanding maneuver pattern in Case~2 requires
failure-aware online reassignment or replanning. We now state this practical
limitation explicitly rather than claiming unconditional fault tolerance.

\begin{figure*}[!t]
  \centering
  \includegraphics[width=\textwidth]
  {figures/two_defender_failure_mc_boxplots.pdf}
  \caption{Terminal-metric distributions under two randomly unavailable
  interceptors: (a) temporal-coordination error, (b) terminal lateral
  acceleration, (c) terminal miss distance, and (d) coordinated interception
  duration. Continuous metrics are evaluated over successful interception
  trials only.}
  \label{fig:two_interceptor_failure_mc}
\end{figure*}

% Important internal note:
% The audited stable branch does not contain an independent friend-to-friend
% communication-delay or packet-loss/hold-last-message implementation.
% Do not claim that the present experiment validates packet loss until that
% mechanism is implemented and evaluated with the same frozen-policy protocol.
